\begin{abstract}
Although large-scale unsupervised language models (LMs) acquire extensive world knowledge and demonstrate certain reasoning abilities, exerting precise control over their outputs remains challenging due to the entirely unsupervised approach used in their pretraining. To achieve greater steerability, current approaches involve gathering human judgments that indicate relative preference among model outputs, subsequently fine-tuning the original LM to reflect these preferences—most often using reinforcement learning from human feedback (RLHF). However, RLHF typically involves a multi-step, often unstable process: first, constructing a reward model that encapsulates the human preferences, and next, using reinforcement learning to update the large LM so as to optimize this reward, all the while constraining deviation from the initial model parameters. This work presents a novel parameterization for the reward model in RLHF, which admits an explicit solution for the associated optimal policy. As a result, the classic RLHF problem can be addressed via a straightforward classification loss. We term our method \textit{Direct Preference Optimization} (DPO), and find that it is stable, achieves strong empirical results, and is significantly less computationally demanding. DPO removes the necessity for sampling from the LM during the adaptation phase and minimizes the need for extensive hyperparameter searches. Empirical studies demonstrate that DPO is able to fine-tune LMs to adhere to human preferences as well as, or more effectively than, existing techniques. In particular, DPO surpasses PPO-based RLHF in controlling generation sentiment, and matches or exceeds the quality of generated summaries and single-turn dialogue responses, all while being considerably more straightforward to implement and train.
\end{abstract}

\section{Introduction}
Unsupervised language models (LMs) trained on massive-scale corpora have demonstrated unexpectedly advanced abilities~\citep{chowdhery2022palm, brown2020language, touvron2023llama, bubeck2023sparks}. Despite these strengths, such models are typically exposed to data generated by humans with diverse objectives, priorities, and varying degrees of expertise. Not all of these human-driven tendencies are desirable for models to emulate; for instance, though an AI programming assistant should \textit{recognize} common errors to better correct them, we still prefer the model's generated code to reflect the exceptional—albeit possibly uncommon—high-quality programming evident in some of its training sources. Analogously, a language model should \textit{recognize} prevalent misconceptions, such as those held by half of the population, but it is undesirable for the model to assert those inaccuracies as true half the time when queried. This highlights the critical need to \emph{selectively elicit the preferred outputs and conduct} from the model's extensive repertoire of \textit{acquired knowledge and capabilities} to ensure that AI systems are reliable, effective, and aligned with human intent~\citep{ouyang2022training}. Although most current strategies aim to align LMs with human values via reinforcement learning (RL), we will demonstrate that the RL-driven objective adopted by these approaches can in fact be solved directly through a basic binary cross-entropy loss, thereby streamlining the entire preference optimization framework.

Generally, standard methods imbue language models with target behaviors through carefully constructed datasets of human preferences that encode what is considered safe and useful. This stage of learning from preferences follows a foundational phase of large-scale unsupervised pre-training using extensive textual data. While supervised fine-tuning on exemplary human demonstrations of desirable outputs offers a straightforward approach to preference learning, methods based on reinforcement learning from human (or artificial) feedback (RLHF/RLAIF; \citep{christiano2017deep,bai2022constitutional}) have proven most effective. RLHF involves training a reward model based on annotated human preference data, followed by using RL to optimize the language model such that it generates outputs with higher predicted reward, all the while limiting deviation from the initial model distribution. Despite the notable advancements in conversation and programming brought about by RLHF, this approach introduces considerable complexity relative to standard supervised learning. It necessitates training several distinct models and repeatedly sampling from the language model policy during optimization, thereby resulting in considerable computational overhead.

This work introduces a method for aligning language models with human preferences that bypasses both explicit reward modeling and reinforcement learning. We introduce \textit{{Direct Preference Optimization} (DPO)}, a training procedure that implicitly targets the same optimization criterion as traditional RLHF techniques—maximizing rewards subject to a KL-divergence penalty—but features a more straightforward design and implementation. Conceptually, DPO updates function by increasing the log odds in favor of preferred over non-preferred model outputs. Importantly, the method integrates a dynamically computed, per-sample importance weight, which mitigates the model collapse issues seen when employing a simplistic probability ratio approach. As with prior frameworks, DPO makes use of a formal preference model (for example, the Bradley-Terry model; \cite{bradley1952rankanalysis}) to quantify how well a reward captures observed preference behavior. Differing from established pipelines, which train a reward model using preference losses derived from these models and subsequently train a policy against the learned reward, DPO reformulates the preference loss to directly depend on the policy, via a reparameterization. This allows DPO to exploit human-annotated preference datasets to optimize the policy through a standard binary cross-entropy loss, thereby achieving the optimal response policy for a latent reward shaped by the preferences.

Our primary contribution is the Direct Preference Optimization algorithm, a straightforward preference-based approach to training language models that does not rely on reinforcement learning. Empirically, we demonstrate that DPO matches or surpasses the performance of prevailing approaches—including PPO-based RLHF—in domains such as sentiment control, summarization, and conversational tasks, evaluated on models as large as 6B parameters.

\section{Related Work}

Progressively larger self-supervised language models have demonstrated the ability to solve certain tasks without any training examples (zero-shot; \citep{radford2019language}) or with only a handful of provided instances (few-shot prompts; \citep{gpt3,megatron,chowdhery2022palm}). Nevertheless, substantial improvements in downstream performance and user intent alignment are observed when these models are further adapted through fine-tuning on curated datasets composed of explicit instructions paired with human-generated completions \citep{mishra-etal-2022-cross,sanh2022multitask,chung2022scaling,thoppilan2022lamda}. This approach, termed `instruction-tuning,' equips large language models to handle instructions that go beyond those seen during fine-tuning, generally enhancing their practical applicability \citep{chung2022scaling}. 

Although instruction tuning is effective, collecting \textit{relative} human judgments—comparisons of response quality—tends to be more feasible than gathering high-quality expert demonstrations. Consequently, later studies have leveraged preference-based datasets for LLM fine-tuning, yielding improvements in a variety of domains such as translation \citep{kreutzer-etal-2018-reliability}, summarization \citep{stiennon2022learning,ziegler2020finetuning}, story generation \citep{ziegler2020finetuning}, and instruction compliance \citep{ouyang2022training,ramamurthy2023is}. These strategies typically start by learning a neural network-based reward function, designed to align with observed preferences under models like the Bradley-Terry framework \citep{bradley1952rankanalysis}. The language model is then fine-tuned to maximize these reward estimates, frequently employing reinforcement learning methods such as REINFORCE \citep{williams1992reinforce}, proximal policy optimization (PPO; \cite{schulman2017proximal}), or their extensions \citep{ramamurthy2023is}.

Another closely associated research avenue employs LLMs already fine-tuned on human feedback to produce synthetic datasets targeting specific criteria, like safety or non-harmfulness \citep{bai2022constitutional}, relying solely on human-provided rubrics as weak supervision for annotation. Altogether, these techniques unify efforts from two research traditions: the application of reinforcement learning to language model training for diverse purposes~\citep{Ranzato2015SequenceLT,paulus2018a,wu2018learning}, and the general study of preference-based learning from human feedback \citep{christiano2017deep,kupcsik2018learning}. Despite the intuitive advantages of using relative human preferences, the process of fine-tuning large-scale language models with reinforcement learning presents significant practical hurdles. The current work offers a principled, theory-driven method to optimize for relative preferences without requiring reinforcement learning.

Beyond language-based tasks, the problem of learning policies from preferences has been explored within both bandit and reinforcement learning frameworks, resulting in a variety of proposed methodologies. The contextual dueling bandit (CDB) paradigm \citep{yue2012karmed,dudik2015contextual} extends contextual bandits by using preference feedback—rankings or pairwise comparisons between actions—instead of numeric rewards. In this framework, where explicit reward signals are absent, the notion of optimality is replaced by the \textit{von Neumann winner}: a policy that is expected to outperform any alternative policy at least half the time \citep{dudik2015contextual}. A notable distinction is that CDB approaches generally operate in an online regime with immediate access to preference information, whereas learning from human preferences often relies on a predetermined, static dataset of action pairs annotated with preferences \citep{yan2022human}. Analogously, in \textit{preference-based RL} (PbRL), policies are trained using binary preferences inferred from an unknown 'scoring' function rather than direct reward signals \citep{BusaFekete2014,ruiz2023dueling}. There exists a spectrum of PbRL algorithms, including those that are capable of leveraging off-policy preference data, but these methods typically involve first estimating the underlying reward or scoring function before proceeding to optimize the policy against it \citep{jain2013learning,BusaFekete2014,christiano2017deep,sadigh2017active,kupcsik2018learning}. In contrast, we propose a unified, single-phase approach where policy parameters are directly adjusted to fulfill observed preferences.

\section{Preliminaries}\label{section:prelims}

We outline the RLHF workflow as introduced by \citeauthor{ziegler2020finetuning}, and subsequently elaborated in works such as \citep{stiennon2022learning, bai2022training, ouyang2022training}. This pipeline commonly consists of three major stages: (1) supervised fine-tuning (SFT), (2) preference-driven reward model training, and (3) reinforcement learning-based policy improvement.

\textbf{SFT}: The first step in RLHF is to further fine-tune a pretrained language model on carefully curated examples related to the end-use application (e.g., conversational agents, text summarization). This results in a supervised model denoted by $\pi^\text{SFT}$. 

\textbf{Reward Modelling Phase}: Next, the SFT model receives an input $x$ and produces response pairs $(y_1, y_2)\sim \pi^\text{SFT}(y \mid x)$. Human annotators then assess these outputs, selecting the preferred completion (denoted $y_w\succ y_l \mid x$, where $y_w$ and $y_l$ correspond to the selected and non-selected responses, respectively). The underlying rationale for these selections is presumed to be governed by an implicit reward function $r^*(y, x)$, which is inaccessible to the modeller. Several models have been developed for translating human preferences into a quantitative form, with the Bradley-Terry (BT) model \cite{bradley1952rankanalysis} being among the most prevalent; other ranking systems such as the Plackett-Luce model \citep{plackett1975analysis, luce2012individual} are applicable when multi-way ranking data are available. According to the BT framework, the human-derived preference distribution $p^*$ is defined as follows:

\begin{equation}\label{eq:bradley-terry}
    p^*(y_1\succ y_2 \mid x)=\frac{\exp\left(r^*(x, y_1)\right)}{\exp\left(r^*(x, y_1)\right) + \exp\left(r^*(x, y_2)\right)}.
\end{equation}

Given access to a fixed set of comparisons $\mathcal{D}=\bigl\{x^{(i)}, y_w^{(i)}, y_l^{(i)}\bigr\}_{i=1}^N$ drawn from $p^*$, we can introduce a parameterized reward function $r_{\phi}(x, y)$ and estimate its parameters by maximizing likelihood. By interpreting this setup as a binary classification problem, the resulting objective becomes the negative log-likelihood:

\begin{equation}\label{eq:reward_model}
    \mathcal{L}_R(r_{\phi}, \mathcal{D}) = -\mathbb{E}_{(x, y_w, y_l)\sim \mathcal{D}}\bigl[\log \sigma(r_{\phi}(x, y_w)- r_{\phi}(x, y_l))\bigr]
\end{equation}

where $\sigma$ represents the logistic sigmoid. In large language models, $r_{\phi}(x, y)$ is typically initialized from the SFT model $\pi^\text{SFT}(y \mid x)$, enhanced by a linear output head atop the last transformer block to yield a single scalar reward \cite{ziegler2020finetuning}. To stabilize the estimated rewards and limit their variance, previous work typically enforces normalization so that $\mathbb{E}_{x,y\sim \mathcal{D}}\left[r_\phi(x, y)\right] = 0$ holds for all $x$.

\textbf{RL Fine-Tuning Phase}: In the RL fine-tuning step, this trained reward model guides updates to the language model policy. In line with prior literature~\citep{jaques2017sequence, jaques2020human}, policy optimization is formalized as

\begin{equation}\label{eq:RL}
\max_{\pi_{\theta}}  \mathbb{E}_{x\sim \mathcal{D}, y\sim \pi_{\theta}(y \mid x)}\bigl[r_{\phi}(x, y)\bigr] - \beta\mathbb{D}_{\textrm{KL}}\bigl[\pi_{\theta}(y\mid x)\mid \mid \pi_\text{ref}(y\mid x)\bigr],
\end{equation}

where $\beta$ determines the degree of regularization with respect to the reference policy $\pi_\text{ref}$, typically set as the initial SFT model $\pi^\text{SFT}$. In practice, the policy $\pi_\theta$ is often seeded with $\pi^\text{SFT}$ as well. The inclusion of the KL-divergence constraint helps restrict the policy from straying excessively from distributions where the reward function has validity, and also promotes diversity in model outputs, thereby reducing the risk of mode collapse on high-reward generations. As language modeling involves discrete outputs, this objective is not directly differentiable and is usually tackled with reinforcement learning algorithms. The most widely-adopted method \citep{ziegler2020finetuning, stiennon2022learning, bai2022training, ouyang2022training} constructs the following reward objective: ${r(x, y) = r_{\phi}(x, y) -\beta (\log \pi_{\theta}(y\mid x) - \log \pi_\text{ref}(y\mid x))}$, and maximizes it using PPO \cite{schulman2017proximal}.

\section{Direct Preference Optimization}\label{sec:DPO}

Seeking alternatives to reinforcement learning algorithms for scaling language model fine-tuning, our objective is to develop a straightforward method for optimizing the policy directly from preference data. Contrary to prior RLHF strategies—wherein a reward model is fit and then optimized via RL—our approach exploits a specific parameterization of the reward function, allowing us to derive the associated optimal policy analytically, thus avoiding RL training altogether. The core idea is to employ an explicit mapping between reward functions and their optimal policies, thereby recasting the original loss over reward functions as an equivalent loss over policies. Through this reparameterization, it becomes unnecessary to independently fit a reward model, yet the procedure retains compatibility with established preference modeling frameworks like the Bradley-Terry model. Effectively, this means the policy network simultaneously embodies both the generative language model and the

\textbf{Derivation of the DPO objective.} Beginning with the standard RL objective defined in Eq.~\ref{eq:RL}, parameterized by a general reward function $r$, we follow previous research~\citep{peters2007reinforcement, peng2019advantage, korbak2022reinforcement, go2023aligning} to observe that the KL-regularized reward maximization problem admits a closed-form optimal solution:

\begin{equation}\label{eq:op_policy}
    \pi_r(y\mid x) = \frac{1}{Z(x)}\pi_\text{ref}(y\mid x)\exp\left(\frac{1}{\beta}r(x, y)\right),
\end{equation}

where $Z(x) =\sum_{y}\pi_\text{ref}(y\mid x)\exp\left(\frac{1}{\beta}r(x, y)\right)$ acts as a normalization constant. The detailed derivation can be found in Appendix \ref{app:derivation1}. When replacing the reward function $r$ with its maximum likelihood estimator $r_{\phi}$ for the true but unknown $r^*$, calculating $Z(x)$ remains intractable for practical applications~\citep{korbak2022reinforcement, go2023aligning}. To alleviate this, we rewrite Eq.~\ref{eq:op_policy} by expressing $r(x, y)$ in terms of the optimal policy $\pi_r$, the reference policy $\pi_\text{ref}$, and $Z(\cdot)$. By applying the logarithm to both sides and manipulating algebraically, we obtain:

\begin{equation}\label{eq:main_eq}
    r(x,y) =\beta \log \frac{\pi_r(y\mid x)}{\pi_\text{ref}(y\mid x)} + \beta \log Z(x).
\end{equation}

This formulation can be applied to the true reward $r^*$ as well as its associated optimal policy $\pi^*$. Notably, the Bradley-Terry preference model evaluates the difference in rewards between two output candidates, given as ${p^*(y_1 \succ y_2 \mid x) = \sigma(r^*(x, y_1) - r^*(x, y_2))}$. Substituting the decomposition in Eq.~\ref{eq:main_eq} into Eq.~\ref{eq:bradley-terry} for $r^*(x,y)$, the normalization terms $Z(x)$ cancel out, allowing us to represent preference probabilities solely in terms of the policies $\pi^*$ and $\pi_\text{ref}$. As a result, the optimal RLHF policy $\pi^*$ under the Bradley-Terry framework is governed by:

\begin{equation}\label{eq:objective}
    p^*(y_1\succ y_2 \mid x)=\frac{1}{1 + \exp\left(\beta \log \frac{\pi^*(y_2\mid x)}{\pi_\text{ref}(y_2\mid x)} - \beta \log \frac{\pi^*(y_1\mid x)}{\pi_\text{ref}(y_1\mid x)}\right)}
\end{equation}

The derivation of this result appears in Appendix~\ref{app:derivation2}. While the discussion here utilizes the Bradley-Terry model, analogous constructions hold for more general ranking models such as the Plackett-Luce model~\citep{plackett1975analysis, luce2012individual}, as described in Appendix~\ref{app:plackett_luce_models}.

With this formulation, where preference probabilities depend only on the optimal policy, we can define a maximum likelihood training objective directly for a parameterized policy $\pi_\theta$. Drawing a parallel to the standard reward modeling approach (see Eq.~\ref{eq:reward_model}), the loss for DPO is given by:

\begin{equation}\label{eq:optimum_model}
    \mathcal{L}_\text{DPO}(\pi_{\theta}; \pi_\text{ref}) = -\mathbb{E}_{(x, y_w, y_l)\sim \mathcal{D}}\left[\log \sigma \left(\beta \log \frac{\pi_{\theta}(y_w\mid x)}{\pi_\text{ref}(y_w\mid x)} - \beta \log \frac{\pi_{\theta}(y_l\mid x)}{\pi_\text{ref}(y_l\mid x)}\right)\right].
\end{equation}

In this formulation, we model the implicit reward through an alternative parameterization, ensuring that the resulting optimal policy coincides with $\pi_\theta$. Additionally, because this approach aligns with a reparametrized Bradley-Terry model, it inherits specific theoretical guarantees, such as consistency when appropriate conditions on the preference data distribution are met \cite{bong2022generalized}. Further discussion regarding DPO's theoretical attributes and its relationship with related frameworks is presented in Section~\ref{sec:theory}.

\textbf{Mechanistic Interpretation of the DPO Update:} To gain a concrete understanding of DPO's functioning, it is helpful to examine the gradient of the objective $\mathcal{L}_\text{DPO}$. The parameter gradient is given by:

\begin{multline*}\label{eq:gradient}
    \nabla_\theta \mathcal{L}_\text{DPO}(\pi_\theta;\pi_\text{ref}) = \\ -\beta\mathbb{E}_{(x, y_w, y_l) \sim \mathcal{D}} \bigg[\underbrace{\sigma(\hat{r}_\theta(x, y_l) - \hat{r}_\theta (x, y_w))}_\text{places greater emphasis where reward ordering is incorrect}\bigg[\underbrace{\nabla_\theta\log \pi(y_w \mid x)}_\text{push up $y_w$'s likelihood} - \underbrace{\nabla_\theta\log\pi(y_l \mid x)}_\text{push down $y_l$'s likelihood}\bigg]\bigg],
\end{multline*}

with $\hat{r}_\theta(x, y) = \beta \log \frac{\pi_\theta(y \mid x)}{\pi_\text{ref}(y \mid x)}$ acting as the reward defined implicitly by both the learned policy $\pi_\theta$ and the reference policy $\pi_\text{ref}$ (see Section~\ref{sec:theory} for elaboration). Conceptually, this gradient modifies $\pi_\theta$ by boosting the probabilities assigned to preferred responses $y_w$ and suppressing those for less preferred alternatives $y_l$. The influence of each data point is scaled according to how much the implicit reward function $\hat{r}_\theta$ erroneously prefers the disfavored choice over the preferred one, with the scaling factor $\beta$ capturing the strength of the KL regularization. Empirical results highlight the significance of this reweighting; omitting it leads to detrimental behavior in the model, as demonstrated by the failure cases in Appendix Table~\ref{tab:unlikelihood_generations}.

\textbf{DPO overview.} The typical DPO process proceeds as follows: 1) For each prompt $x$, generate completions $y_1, y_2$ by sampling from $\pi_\text{ref}(\cdot \mid x)$ and assign human preference labels, thereby forming the offline preference dataset $\mathcal{D} = \{x^{(i)}, y_w^{(i)}, y_l^{(i)}\}_{i=1}^N$; 2) train the language model $\pi_\theta$ by minimizing $\mathcal{L}_\text{DPO}$ with respect to the fixed $\pi_\text{ref}$, the dataset $\mathcal{D}$, and a specified value of $\beta$.

Practically, utilizing existing public preference datasets is preferable over generating new samples and collecting human annotations. Because such datasets are derived using $\pi^\text{SFT}$, we set $\pi_\text{ref} = \pi^\text{SFT}$ whenever possible. In scenarios where $\pi^\text{SFT}$ is not provided, we instead set $\pi_\text{ref}$ to maximize the likelihood of preferred responses ${(x, y_w)}$, expressed as ${\pi_\text{ref} = \argmax_{\pi}\mathbb{E}_{x, y_w \sim \mathcal{D}}\left[\log \pi(y_w \mid x)\right]}$. This approach helps to counteract the mismatch between the inaccessible true reference distribution and the $\pi_\text{ref}$ utilized within DPO. Additional technical and hyperparameter details are discussed in Appendix~\ref{app:implementation}.

\section{Theoretical Analysis of DPO}
This section provides a deeper theoretical examination of the DPO approach, articulates its theoretical foundation, and elucidates the benefits of DPO in relation to the limitations encountered by actor-critic methods in RLHF, such as PPO~\cite{schulman2017proximal}.

\label{sec:theory}

\subsection{Your Language Model Is Secretly a Reward Model} 
DPO circumvents the need to fit an explicit reward model and to perform reinforcement learning by leveraging a single maximum likelihood objective. The minimization objective given by Eq. \ref{eq:main_eq} can be interpreted as the Bradley-Terry model with a reward representation $r^*(x, y) = \beta \log\frac{\pi^*_\theta(y \mid x)}{\pi_\text{ref}(y \mid x)}$, where the parametric model $\pi_{\theta}$ is trained analogously to reward model optimization (see Eq. \ref{eq:reward_model}) after an appropriate variable substitution. The following discussion constructs the formal framework supporting this reparameterization, demonstrates that it imposes no additional limitations on the reward models that can be learned, and proves that it enables recovery of the optimal policy exactly. We begin by introducing an equivalence relation among reward functions.

\begin{definition}
Two reward functions $r(x, y)$ and $r'(x, y)$ are considered equivalent if and only if there exists a function $f$ such that $r(x, y)-r'(x, y) = f(x)$.     
\end{definition}

This equivalence is straightforward to verify, and it partitions all reward functions into distinct equivalence classes. We now present two immediate results:

\begin{lemma}\label{lemma:same_prefrence} Within the Plackett-Luce framework, including the specific case of the Bradley-Terry model, any two reward functions belonging to the same equivalence class produce identical distributions over preferences.
\end{lemma}

\begin{lemma}\label{lemma:same_policy}
    Any two reward functions from the same equivalence class will result in the same optimal policy for the associated constrained RL problem.
\end{lemma}

Since the proofs are elementary, they are included in Appendix \ref{app:lemma1}. The first lemma reflects the widely recognized issue of under-specification present in the Plackett-Luce family of preference models \cite{plackett1975analysis}. As a result of this ambiguity, it is common practice to impose extra identifiability restrictions to ensure meaningful maximum likelihood estimates from Eq. \ref{eq:reward_model} \cite{bong2022generalized}. According to the second lemma, all members of a given equivalence class give rise to the same optimal policy, so our ultimate aim is simply to recover some reward function within the optimal equivalence class. The following theorem, proven in Appendix~\ref{app:thm1}, formalizes this property:

\begin{theorem}\label{thm:main}
    Given mild conditions, every equivalence class of rewards consistent with the Plackett-Luce (and specifically Bradley-Terry) models admits the following parameterization: ${r(x, y) = \beta \log \frac{\pi(y\mid x)}{\pi_\text{ref}(y\mid x)}}$ for some choice of model $\pi(y\mid x)$ and for any fixed reference policy $\pi_\text{ref}(y \mid x)$.
\end{theorem}

\begin{sproof}
    Take any reward function $r(x, y)$, which determines an optimal model $\pi_r(y \mid x)$ as described in Eq. \ref{eq:op_policy}. We now demonstrate that there is a reward function in the equivalence class of $r$ that can be expressed through the reparameterization introduced above. To this end, define the transformation $f$ by  
\begin{equation}
    f(r; \pi_\text{ref}, \beta)(x, y) = r(x, y) - \beta\log\sum_{y}\pi_\text{ref}(y\mid x)\exp\left(\frac{1}{\beta}r(x, y)\right)
\end{equation}
Here, $f$ serves to adjust the reward function by subtracting the log-partition function associated with $\pi_r$. Because this normalization term depends solely on the context $x$, $f(r; \pi_\text{ref}, \beta)(x, y) $ remains within the equivalence class of $r(x, y)$. Now, substituting the right side of Eq.~\ref{eq:main_eq} (which is valid for any reward), it follows that $f(r; \pi_\text{ref}, \beta)(x, y) = \beta \log \frac{\pi_r(y\mid x)}{\pi_\text{ref}(y\mid x)}$. This shows that the mapping $f$ generates a representative of $r$'s equivalence class in the specified form, confirming that our reparameterization does not restrict the generality of the reward model.
\end{sproof}

An alternative interpretation of Theorem~\ref{thm:main} is that it pinpoints the precise reward function chosen by the DPO reparameterization within each equivalence class; specifically, it is the reward function that satisfies

\begin{equation}\label{eq:lag_p}
     \sum_{y}\underbrace{\pi_\text{ref}(y\mid x)\exp\left(\frac{1}{\beta}r(x, y)\right)}_{=\pi(y\mid x)\text{, using Thm.~\ref{thm:main} reparam.}} = 1,
\end{equation}

which asserts that $\pi(y\mid x)$ constitutes a legitimate probability distribution—that is, all probabilities are nonnegative and the total probability sums to one. Notably, referencing Eq.~\ref{eq:op_policy}, one can recognize Eq.~\ref{eq:lag_p} as describing the partition function corresponding to the optimal policy determined by $r(x, y)$. The central observation underlying the DPO algorithm is that by enforcing suitable constraints on the inherently underdetermined Plackett-Luce (and particularly Bradley-Terry) preference modeling family, it is possible to maintain the expressivity over reward functions while ensuring the optimal policy in Eq. \ref{eq:op_policy} has an analytically manageable form for every prompt $x$.

\subsection{Instability of Actor-Critic Algorithms} Our framework also facilitates the identification of instabilities present in standard actor-critic algorithms, such as PPO, commonly employed in RLHF setups. Here, we adhere to the RLHF methodology, concentrating on the RL fine-tuning phase detailed in Section \ref{section:prelims}. Connections may be established with the control as inference perspective \cite{levine2018reinforcement}, as applied to the constrained RL formulation provided in \ref{eq:RL}. Suppose we work with a parameterized policy $\pi_{\theta}(y\mid x)$; the objective becomes minimizing $\mathbb{D}_{\text{KL}}[\pi_{\theta}(y|x) \mid \mid \pi^*(y\mid x)]$, where $\pi^*$ represents the optimal policy dictated by Eq. \ref{eq:optimum_model}, based on reward $r_{\phi}(y, x)$. With some algebraic manipulations, the problem can be cast as the following optimization:

\begin{equation}\label{eq:AC}
    \max_{\pi_{\theta}}\mathbb{E}_{\pi_{\theta}(y\mid x)}\left[\underbrace{r_{\phi}(x, y) -\beta\log\sum_{y}\pi_\text{ref}(y\mid x)\exp\left(\frac{1}{\beta}r_{\phi}(x, y)\right)}_{f(r_{\phi}, \pi_\text{ref}, \beta)} - \underbrace{\beta\log\frac{\pi_{\theta}(y\mid x)}{\pi_\text{ref}(y\mid x)}}_{\text{KL}}\right]
\end{equation}

The objective considered here is identical to that targeted in previous studies 
\citep{ziegler2020finetuning, stiennon2022learning, bai2022training, ouyang2022training}, which utilize a DPO-equivalent reward for the reward family $r_{\phi}$. In this context, the normalization component within $f(r_{\phi}, \pi_\text{ref}, \beta)$ can be viewed as the soft value function corresponding to the reference policy $\pi_\text{ref}$. Although this normalization does not influence the optimal solution itself, omitting it could substantially increase the variance of policy gradients, leading to less stable optimization. To account for normalization, one might employ a trainable value function, but this route often encounters optimization difficulties. Previous work often circumvents this by normalizing rewards using a human completion baseline, effectively estimating the normalizing term via a single-sample Monte Carlo method. In contrast, the DPO reparameterization produces a reward function inherently free from such baseline requirements.

\section{Experiments}
This section presents an empirical analysis of DPO for policy training based solely on preference data. We start by investigating a controlled text-generation task to address the following: how does DPO manage the trade-off between maximizing reward and minimizing KL-divergence to the reference policy in comparison to standard preference optimization methods, such as PPO? Subsequently, we assess the effectiveness of DPO on larger-scale models and on more challenging RLHF benchmarks, like summarization and dialogue. Our results indicate that, even with minimal hyperparameter adjustment, DPO generally matches or outperforms established approaches—including RLHF with PPO and selecting the best out of $N$ sampled trajectories as evaluated by a trained reward model. We detail our experimental methodology below, and further specifics can be found in Appendix~\ref{app:exp_details}.

\textbf{Tasks.} Our study examines three distinct open-ended text generation challenges. Across all tasks, the algorithms are trained to derive a policy based on a dataset of human preferences denoted as $\mathcal{D}=\bigl\{x^{(i)}, y_w^{(i)}, y_l^{(i)}\bigr\}_{i=1}^N$. In the \textbf{controlled sentiment generation} task, $x$ consists of an initial segment from a movie review extracted from the IMDb dataset \cite{maas-EtAl:2011:ACL-HLT2011}. The goal for the policy is to generate $y$ with a positive affect. To systematically assess performance, preference pairs of generated outputs are automatically constructed with the aid of a pre-trained sentiment classifier; pairs are selected such that $p(\text{positive}\mid x,y_w)>p(\text{positive}\mid x,y_l)$. For SFT, GPT-2-large is further trained to convergence on training data taken from the IMDb review dataset (additional implementation details can be found in App~\ref{app:sentiment_details}). In the \textbf{summarization} task, $x$ corresponds to a Reddit forum post, and the aim is to produce a summary $y$ that captures its principal ideas. As in earlier studies, we make use of the Reddit TL;DR summarization dataset \citep{volske-etal-2017-tl} in conjunction with a set of human preferences assembled by \citeauthor{stiennon2022learning}. The SFT model used here is fine-tuned on human-composed summaries of forum posts\footnote{\url{https://huggingface.co/CarperAI/openai_summarize_tldr_sft}} leveraging the TRLX \citep{leandro_von_werra_2023_7790115} RLHF toolkit. The dataset containing human judgments was created by \citeauthor{stiennon2022learning} based on outputs from another SFT model with similar training parameters. Lastly, the \textbf{single-turn dialogue} task sets $x$ as a user-issued query, which might span a range of topics, from astrophysics inquiries to advice requests. Here, the policy is expected to generate $y$, a helpful and engaging reply. We utilize the Anthropic Helpful and Harmless dialogue dataset \citep{bai2022training}, which consists of 170,000 interactions between people and an automated agent. Each interaction concludes with two machine-generated replies and a label indicating the human-chosen preferred response. Since no pre-existing SFT model is available for this domain, we construct one by fine-tuning an off-the-shelf language model exclusively on the completions marked as preferred.

\textbf{Evaluation.} Our experimental framework employs two primary evaluation strategies. For controlled sentiment generation, we assess the capability of each algorithm to optimize the constrained reward maximization objective by charting the achievable trade-off between reward and KL-divergence from the reference policy; such a trade-off frontier can be directly computed owing to our access to the true reward signal (as measured by a sentiment classifier). In contrast, since real-world settings lack access to the ground-truth reward function, we instead determine algorithm performance by measuring the \textit{win rate} over a baseline policy, with GPT-4 serving as an automated proxy for human judgments. Specifically, we utilize GPT-4 to compare systems on summary quality and dialogue helpfulness within summarization and single-turn dialogue tasks, respectively. For summarization, the baseline corresponds to reference summaries drawn from the evaluation set, while in dialogue the baseline is set as the annotator-preferred responses. Given evidence that large language models (LMs) are more reliable evaluators than traditional metrics \citep{Chen2023ExploringTU}, we further substantiate our reliance on GPT-4 via a dedicated human evaluation study discussed in Sec.~\ref{sec:human-judgments}. This study reveals that GPT-4's assessments are highly correlated with human judgments, with the degree of agreement between humans and GPT-4 often matching or surpassing inter-annotator consistency.

\textbf{Methods.} Beyond DPO, our comparisons include a variety of established techniques for aligning language models with human preferences. As baseline approaches, we test zero-shot prompting of \textbf{GPT-J} \citep{gpt-j} in the summarization domain, and two-shot prompting of \textbf{Pythia-2.8B} \citep{biderman2023pythia} within dialogue. We also evaluate the \textbf{SFT} model, as well as \textbf{Preferred-FT}—a supervised model fine-tuned on preferred completions $y_w$ generated either by the SFT model (in both controlled sentiment and summarization tasks) or a generic language model (for single-turn dialogue). Another approach assessed is the pseudo-supervised \textbf{Unlikelihood}~\citep{welleck2019neural} objective, which encourages the policy to increase the probability assigned to $y_w$ while simultaneously \textit{decreasing} probability on $y_l$, utilizing a tunable coefficient $\alpha\in[0,1]$ on the 'unlikelihood' component. We also include \textbf{PPO} \citep{schulman2017proximal}, trained using a reward model derived from preference data, and \textbf{PPO-GT}, an oracle variant utilizing direct access to the ground-truth reward signal available in sentiment-controlled tasks. In our sentiment generation experiments, two PPO-GT variants are considered: an off-the-shelf implementation \cite{leandro_von_werra_2023_7790115} and a customized version incorporating reward normalization and further hyperparameter refinement to enhance results (these enhancements are also applied when using PPO with learned rewards). Lastly, the \textbf{Best of $N$} baseline involves generating $N$ samples from either the SFT model (or Preferred-FT for dialogue), selecting the top candidate as ranked by a learned reward model. Although this decouples reward modeling from policy optimization—yielding strong performance—the approach is computationally infeasible at inference for even moderate values of $N$, as it necessitates sampling $N$ completions for every test query.

\subsection{How well can DPO optimize the RLHF objective?}

In standard RLHF algorithms, the objective involves maximizing the reward while incorporating a KL penalty to limit how much the learned policy diverges from the reference policy. Thus, in evaluating and comparing different algorithms, it is essential to consider not just the reward obtained but also the KL divergence; an increase in reward is not preferable if it comes at the cost of a significant rise in KL. The reward-KL frontiers for various algorithms in the sentiment task are visualized in Figure~\ref{fig:frontier-tldr-main}. For each method, we conduct several training runs, each varying the hyperparameter controlling policy conservativeness (target KL values of $\{3,6,9,12\}$ for PPO, $\beta$ values of $\{0.05,0.1,1,5\}$, $\alpha$ values of $\{0.05,0.1,0.5,1\}$ for unlikelihood, and different random seeds for preferred-FT), totaling 22 runs. At intervals of 100 training steps up to convergence, we assess each policy using a set of test prompts, calculating both the mean reward using the actual reward function and the mean sequence-level KL\footnote{Calculated as the sum of KL divergences at each timestep.} between the current and reference policy, denoted as $\text{KL}\left(\pi\mid \mid \pi_\text{ref}\right)$. The analysis demonstrates that DPO outperforms others in producing an efficient reward-KL frontier: it achieves higher reward levels while keeping the KL low. This superiority is striking for several reasons. First, despite DPO and PPO targeting the same objective, DPO attains substantially greater efficiency; its tradeoff between reward and KL is superior to that of PPO. Moreover, DPO establishes a better reward-KL frontier compared to PPO, \emph{even in the case where PPO is supplied with the actual reward values} (PPO-GT).

\subsection{Can DPO scale to real preference datasets?}
\label{sec:dpo-real-datasets}

We proceed by assessing how DPO performs when fine-tuned on real-world preference datasets in tasks like summarization and single-turn dialogue. In the context of summarization, traditional automatic metrics such as ROUGE have been shown to correlate weakly with human judgments~\citep{stiennon2022learning}, and earlier research demonstrated that LMs optimized via PPO using human feedback tend to generate more desirable summaries. To compare various techniques, we generate outputs from models on the test portion of the TL;DR summarization corpus, then measure the average frequency with which these outputs are favored over references in the test set. Sampling is conducted at multiple temperatures ranging from 0.0 to 1.0; results are visualized in Figure~\ref{fig:frontier-tldr-main} (right). For all methods—DPO, PPO, and Preferred-FT—the identical GPT-J SFT model was used as a base for fine-tuning\footnote{\url{https://huggingface.co/CarperAI/openai_summarize_tldr_sft}}. Experimental findings indicate that DPO attains a win rate near 61\% at a sampling temperature of 0.0, surpassing PPO, which achieves approximately 57\% under optimal conditions at the same temperature. DPO also obtains a higher best-case win rate when compared with the strongest $N$ baseline. It should be mentioned that DPO's $\beta$ parameter was not systematically tuned in these experiments, suggesting that actual DPO performance may be even greater. Additionally, DPO exhibits significantly greater resilience to changes in sampling temperature relative to PPO, whose effectiveness diminishes to that of the baseline GPT-J model at elevated temperatures. No meaningful improvement is observed with Preferred-FT relative to the SFT baseline. Finally, a direct comparison between DPO and PPO through human evaluations, discussed in Section~\ref{sec:human-judgments}, shows that human annotators favored DPO samples at temperature 0.25 over PPO samples at the same temperature 58\% of the time.

For single-turn dialogue tasks, we assess various approaches using the portion of the Anthropic HH dataset \citep{bai2022training} test set comprising single human-assistant exchanges. For evaluations involving GPT-4, the human-preferred responses from the test set serve as the gold standard, enabling calculation of each method’s win rate. In the absence of a canonical SFT model for this benchmark, our pipeline initiates with a pre-trained Pythia-2.8B. We then employ Preferred-FT, finetuning a reference model on the selected completions to ensure alignment with the model’s output distribution, and subsequently apply DPO for further training. The comparative analysis also includes the Best of 128 Preferred-FT completions (as we observe, via Appendix Figure~\ref{fig:best-of-n}, that performance saturates around 128 candidates), and a two-shot prompting configuration with the Pythia-2.8B base model. Across optimal temperature settings, DPO matches or surpasses these baselines. Additionally, we benchmark against an RLHF model utilizing PPO, trained on the Anthropic HH dataset \footnote{\url{https://huggingface.co/reciprocate/ppo_hh_pythia-6B}} following established methodologies \footnote{\url{https://github.com/CarperAI/trlx/tree/main/examples/hh}}. However, in our experiments, no prompt or sampling temperature consistently outperforms the raw Pythia-2.8B baseline with this PPO model. Drawing from outcomes in TL;DR, and noting both Best of 128 and PPO are directed by identical reward objectives, we treat Best of 128 as an approximate surrogate for PPO-level efficacy. In sum, DPO uniquely achieves computationally efficient improvements over human-preferred completions in the Anthropic HH dataset, while rivaling or surpassing the resource-intensive Best of 128 approach. As further illustrated in Figure~\ref{fig:dialogue-main}, DPO rapidly attains peak performance.

\subsection{Generalization to a new input distribution}

To more thoroughly assess how PPO and DPO perform when faced with distribution shifts, we test the policies obtained from our Reddit TL;DR summarization experiment on a distinct dataset—specifically, the test split of CNN/DailyMail containing news articles \citep{nallapati-etal-2016-abstractive}. For these evaluations, we retain the optimal sampling temperatures identified from the TL;DR setup (0 and 0.25). The corresponding outcomes are summarized in Table~\ref{tab:ood}. We measured the win rate of GPT-4 compared to the ground-truth news article summaries, applying the same GPT-4 (C) evaluation prompt as in the Reddit TL;DR experiments, with the term ``forum post'' replaced by ``news article.'' On this alternative data distribution, DPO maintains a notable lead over the PPO policy. These findings offer preliminary support that DPO-trained policies are capable of generalizing to out-of-distribution data at least as effectively as PPO policies, even though DPO does not rely on the supplementary, unlabeled Reddit TL;DR prompts available to PPO.

\subsection{Validating GPT-4 judgments with human judgments}
\label{sec:human-judgments}
A human evaluation was conducted to establish the credibility of GPT-4's assessments by leveraging outcomes from the TL;DR summarization experiment and evaluating two distinct GPT-4 prompts. The \textbf{GPT-4 (S)} prompt (simple) requests a decision on which summary most effectively captures the essential content of the post. The \textbf{GPT-4 (C)} prompt (concise) asks not only about informativeness but also conciseness, as we observe that the \textbf{GPT-4 (S)} configuration often favors summaries that are longer and more repetitive than human raters typically prefer. Full prompt wording is available in Appendix~\ref{app:prompts}. For our analysis, we selected three systems representing a range of output qualities: DPO at temperature 0.25 (highest-performing), PPO at temperature 1.0 (lowest-performing), and SFT at temperature 0.25 (mid-performing), each compared to greedily-sampled PPO (its best temperature setting). This sampling aimed to capture a spectrum of summary quality. Our findings indicate that, under both prompts, GPT-4’s preferences align with human judgments at approximately the same rate as inter-human agreement, implying that GPT-4 serves as an effective proxy for human evaluations. Due to a limited pool of raters, multiple human annotations were only gathered for the DPO and PPO-1 comparisons. Generally, the \textbf{GPT-4 (C)} prompt yields win rates that more closely mirror human preferences, leading us to rely on this prompt for the principal results reported in Section~\ref{sec:dpo-real-datasets}. Further information on the human study, including the rater-facing web interface and participant roster, can be found in Appendix~\ref{app:human-study}.

\section{Discussion}
Preference-based learning presents an effective and scalable approach for developing language models that are both competent and aligned with user values. In this work, we proposed DPO, a straightforward methodology that enables language models to learn from preference data without the need for reinforcement learning. Instead of forcing preference learning into the framework of classical RL to utilize generic RL methods, DPO leverages a connection between the space of language model policies and reward functions. This allows direct optimization of human preferences through a simple cross-entropy loss, omitting the necessity for reinforcement learning while preserving general applicability. DPO matches or outperforms standard RLHF algorithms—such as those using PPO—requiring almost no hyperparameter adjustment. As a result, DPO significantly lowers the practical difficulty of fine-tuning language models from human preference data.

\textbf{Limitations \& Future Work.} Our findings suggest several avenues for further investigation. For instance, how does the DPO-trained policy behave on data distributions that diverge from the training set compared to policies trained via explicit reward modeling? Preliminary evidence points to comparable generalization between DPO and PPO-based approaches, but a systematic exploration is warranted. Another open question is whether leveraging self-labeling with DPO can capitalize on unlabeled prompts as effectively. It is also worth examining the nature of reward over-optimization within DPO; the modest reduction in performance observed in Figure~\ref{fig:dialogue-main}-right might be symptomatic of this effect. Additionally, although our experiments extend to models with up to 6B parameters, future work should investigate the scalability of DPO to models of much greater size. Regarding model assessments, we noticed that GPT-4-derived win rates are sensitive to the choice of prompt; subsequent research could identify optimal strategies to obtain reliable evaluations from automated systems. Ultimately, DPO's applicability may extend far beyond language model alignment, offering promising opportunities for training generative models across diverse domains.